%! Setting Up a Test for the Difference of Two Population Proportions — Unit 6, Lesson 6.10 — Homework (Answer Key)
\documentclass[10pt]{article}
\usepackage{apstats-article}
\usepackage{apstats-key}   % pulls in -boxes; do NOT also load -boxes
\usepackage{grading}
\gradingassignment{apstats-u06-l10-hw}
\gradingcourse{advanced_statistics}
\newcommand{\TallMath}[1]{$\displaystyle #1\rule[-1.4em]{0pt}{3.2em}$}

\begin{document}

\pageheader{Unit 6, Lesson 6.10}{Homework --- Answer Key}

\vspace{0.15in}

\begin{enumerate}

\item \textbf{College Club Membership.} First-gen: 48/160. Non-first-gen: 72/200.

\begin{enumerate}[label=\alph*.]
  \item \ansline{$H_0: p_1=p_2$ (club membership rates are equal for both groups); $H_a: p_1\ne p_2$ (rates differ). Two-sided.}\gradedkeyline{\gradedblank[id=hw1a, type=rubric, points=2, display={$H_0\colon p_1 = p_2$; $H_a\colon p_1 \ne p_2$ (true club-joining proportions, first-gen vs.\ not).}]{\linewidth+6pt}{\gradedlinesheight{1}}}
  \item $\hat p_1 = 48/160 = \ans{0.300}$; \quad $\hat p_2 = 72/200 = \ans{0.360}$; \quad
        $\hat p_c = (48+72)/(160+200) = 120/360 = \ans{0.333}$\gradedkey{$\hat p_1 = $~\gradedblank[id=hw1b1, type=short, points=1, answer=0.300, check=numeric, tolerance=0.001]{0.7in}{0.17in}\quad $\hat p_2 = $~\gradedblank[id=hw1b2, type=short, points=1, answer=0.360, check=numeric, tolerance=0.001]{0.7in}{0.17in}\quad $\hat p_c = $~\gradedblank[id=hw1b3, type=short, points=1, answer=0.333, check=numeric, tolerance=0.001]{0.7in}{0.17in}}
  \item \textbf{Random:} \ansline{Both samples randomly selected --- stated.}\\
        \textbf{Large Counts ($\hat p_c=1/3$):} $160(1/3)=53.3$, $160(2/3)=106.7$, $200(1/3)=66.7$, $200(2/3)=133.3$ --- \ans{all $\ge10$. \checkmark}\\
        \textbf{10\%:} \ansline{Assumed both samples $\le10\%$ of all students in each group.}\gradedkeyline{\gradedblank[id=hw1c, type=rubric, points=3, display={Random samples; with $\hat p_c = 1/3$: $53.3, 106.7, 66.7, 133.3$ all $\ge 10$; each sample $\le 10\%$ of its group.}]{\linewidth+6pt}{\gradedlinesheight{2}}}
\end{enumerate}

\item \textbf{Defect Rate.} Old: 36/400. New: 20/400.

\begin{enumerate}[label=\alph*.]
  \item \ansline{$p_O$ = true defect rate on old line; $p_N$ = true defect rate on new line.\\
        $H_0: p_O=p_N$; $H_a: p_O > p_N$ (old line has higher defect rate). One-sided right.}\gradedkeyline{\gradedblank[id=hw2a, type=rubric, points=2, display={$p_O, p_N$: true defect rates of the old / new line. $H_0\colon p_O = p_N$; $H_a\colon p_O > p_N$.}]{\linewidth+6pt}{\gradedlinesheight{1}}}
  \item $\hat p_c = (36+20)/(400+400) = 56/800 = \ans{0.070}$\gradedkey{$\hat p_c = $~\gradedblank[id=hw2b, type=short, points=1, answer=0.070, check=numeric, tolerance=0.001]{0.8in}{0.17in}}
  \item \textbf{Random:} \ansline{Units randomly selected from each production line --- assumed.}\\
        \textbf{Large Counts ($\hat p_c=0.07$):} $400(0.07)=28$, $400(0.93)=372$ --- \ans{all $\ge10$ for both lines. \checkmark}\\
        \textbf{10\%:} \ansline{Assumed both samples $\le10\%$ of all units from each line.}\gradedkeyline{\gradedblank[id=hw2c, type=rubric, points=3, display={Random samples from each line; with $\hat p_c = 0.07$: $28, 372$ (each line) all $\ge 10$; each sample $\le 10\%$ of its line's units.}]{\linewidth+6pt}{\gradedlinesheight{2}}}
\end{enumerate}

\end{enumerate}

\begin{reflectionbox}
\textbf{Preview:} Lesson 6.11 completes the test: compute the $z$-statistic using $\hat p_c$,
find the $p$-value, and write a full conclusion.
\end{reflectionbox}

\end{document}
